\documentclass[11pt,a4paper]{article}

\usepackage[numbers,sort&compress]{natbib}
\usepackage{amsmath,amssymb}
\usepackage{graphicx}
\usepackage[export]{adjustbox}
\usepackage{capt-of}
\usepackage{booktabs}
\usepackage{array}
\usepackage{tabularx}
\usepackage{multirow}
\usepackage[margin=2.2cm]{geometry}
\usepackage[colorlinks=true,linkcolor=black,citecolor=black,urlcolor=blue]{hyperref}

\graphicspath{{figures/}{./figures/}{./}}
\DeclareGraphicsExtensions{.pdf,.png,.jpg,.jpeg}

\newcolumntype{Y}{>{\raggedright\arraybackslash}X}
\newcolumntype{Z}[1]{>{\centering\arraybackslash}p{#1}}

\renewcommand{\thesection}{S\arabic{section}}
\renewcommand{\thesubsection}{S\arabic{section}.\arabic{subsection}}
\renewcommand{\thefigure}{S\arabic{figure}}
\renewcommand{\thetable}{S\arabic{table}}
\renewcommand{\theequation}{S\arabic{equation}}

\newcommand{\paperfig}[5][]{%
  \begin{center}
  \begin{minipage}{\linewidth}
  \centering
  \includegraphics[
    width=#2,
    max height=0.72\textheight,
    keepaspectratio,
    #1
  ]{#3}%
  \vspace{-0.35em}
  \captionof{figure}{#4}
  \label{#5}
  \end{minipage}
  \vspace{-0.45em}
  \end{center}
}

\newenvironment{suppTableBlock}[2]{%
  \begin{center}
  \begin{minipage}{\linewidth}
  \captionof{table}{#1}
  \label{#2}
  \centering
}{%
  \end{minipage}
  \vspace{0.35em}
  \end{center}
}

\title{\textbf{Supplementary Material for}\\
Joint Model Selection and Adaptation for Few-Shot Digital Twin
Instantiation under Deployment Limits}
\author{}
\date{}

\begin{document}
\maketitle

This supplementary material provides implementation, analytical, and
diagnostic details referenced in the main manuscript. It includes the
operation-count and parameter-count rules, candidate-bank records,
relative-margin calibration, additional analytical results, target-side baseline
settings, diagnostic comparisons, hardware profiling, configuration
analysis, external-transfer results, and synthetic-benchmark generation
details.

The source, architecture-screening, margin-calibration, diagnostic, and
held-out target pools use disjoint center identities.

\section{Detailed Operation-Count and Parameter-Count Rules}
\label{supp:counting}

The estimated operation count is computed for one forward pass with batch
size one. Additions and multiplications in the main linear, convolutional,
recurrent-matrix, residual-projection, and output-layer operations are
counted. Bias additions, nonlinear activations, dropout operations, and GRU
elementwise operations are excluded. The same deterministic rule is applied
to every candidate.

For an instantiated architecture $a_q$, the parameter count is the exact
number of elements in all trainable parameter tensors after the target input
dimension and prediction horizon are fixed. Both counts depend on the target
case only through its instantiated input and output dimensions. Parameter
adaptation changes neither count.

The two indicators are used as hard deployment limits. They are not added
to the prediction loss and are not changed after target adaptation.
They are deterministic model-level indicators rather than direct guarantees
of hardware-specific latency or memory use.
Their relation to measured inference latency and serialized model size is
evaluated in Supplementary Section~S5.

\section{Candidate-Bank Records and Relative-Margin Calibration}
\label{supp:bank-calibration}

\subsection{Candidate-bank construction and frozen records}

The initial architecture space contains 66 indexed configurations: 18 MLPs,
36 TCNs, and 12 GRUs. Configuration A57, a one-layer GRU with 32 hidden
units, is fixed as the reference architecture before screening.

An alternative configuration is retained when it records at least two
check-oracle wins over A57, or at least two validation-selected positive check
wins over A57.
These cases belong to the architecture-screening pool and are disjoint from
the margin-calibration and held-out target pools.
The retained list contains three MLP and three GRU configurations.
No TCN configuration satisfies this rule.

Two retained one-layer GRU-32 configurations have different indexed
settings, but recurrent dropout is inactive for a one-layer GRU. They
therefore represent the same executable architecture. The six retained
configurations correspond to five executable architectures.

The final bank contains the frozen pooled-source PT+FT initialization for
the reference, one source-trained initialization for each of the five other
retained indexed configurations, and a second A57 initialization used as a
paired-initialization control. The six retained configurations therefore
produce seven candidates. The bank and reference candidate are frozen
before relative-margin calibration and held-out evaluation.

\subsection{Relative-margin calibration formulation}

Let
$\mathcal{T}=\{0.05,0.075,0.10,0.125,0.15,0.20\}$
denote the preset margin grid. For a relative margin $\tau$, let
$\widehat R_{\mathrm{harm}}(\tau)$ denote the observed harmful-selection
rate. Let $\widehat G_{\mathrm{ref}}(\tau)$ and
$\widehat G_{\mathrm{pair}}(\tau)$ denote the mean check-loss gains over the
adapted reference and the paired-initialization control. Their lower 95\%
confidence bounds are
$\underline G_{\mathrm{ref}}(\tau)$ and
$\underline G_{\mathrm{pair}}(\tau)$.

The eligible margins form
\begin{equation}
\begin{aligned}
\mathcal{T}_{\mathrm{elig}}
=
\bigl\{
\tau\in\mathcal{T}
\;\bigm|\;&
\widehat R_{\mathrm{harm}}(\tau)\leq\rho,
\\
&
\widehat G_{\mathrm{ref}}(\tau)\geq\gamma_{\mathrm{ref}},
\quad
\underline G_{\mathrm{ref}}(\tau)>0,
\\
&
\widehat G_{\mathrm{pair}}(\tau)\geq\gamma_{\mathrm{pair}},
\quad
\underline G_{\mathrm{pair}}(\tau)>0
\bigr\}.
\end{aligned}
\label{eq:supp-margin-calibration}
\end{equation}

The calibration uses
$\rho=0.05$ and
$\gamma_{\mathrm{ref}}=\gamma_{\mathrm{pair}}=0.03$.
When $\mathcal{T}_{\mathrm{elig}}$ is nonempty, the calibrated margin is
\begin{equation}
\tau^{*}=\min \mathcal{T}_{\mathrm{elig}}.
\end{equation}
The smallest eligible value in the main margin-calibration setting is
$\tau^{*}=0.10$.

Architecture screening and candidate-bank construction are completed before
margin calibration.
The candidate bank and reference candidate are fixed when the margin grid is
evaluated.

If no margin is eligible in a new domain, no domain-specific calibrated
margin is claimed. In calibrated deployment, the adapted reference is
retained and the domain is marked as unsupported for margin calibration.
A previously frozen margin can still be evaluated without recalibration as
a transfer audit, but it is not treated as a calibrated deployment rule for
the new domain. The released
calibration CSV contains the complete six-margin grid, harmful-selection
rates, mean gains, and confidence bounds.

\section{Additional Analytical Details}
\label{supp:analysis}

\subsection{Ideal target decision}

Let $R_c(q)$ denote the expected loss of an adapted candidate on future
observations from target center $c$. If this quantity were available, the
ideal feasible candidate would satisfy
\begin{equation}
q_c^{\dagger}
\in
\underset{q\in\mathcal{Q}_{c}^{\mathrm{fea}}}
{\arg\min}\;
R_c(q).
\label{eq:supp-ideal-target}
\end{equation}

The quantity $R_c(q)$ is not available during model instantiation. MSA-DTI
therefore uses the observed validation loss $J_c(q)$ together with the
relative-margin rule in the main manuscript.

\subsection{Conditional expected-loss result}

Assume that validation loss differs from expected future loss by at most
$\epsilon_c\geq0$ for every feasible candidate:
\begin{equation}
\left|
R_c(q)-J_c(q)
\right|
\leq
\epsilon_c,
\qquad
\forall q\in\mathcal{Q}_{c}^{\mathrm{fea}}.
\label{eq:supp-validation-error}
\end{equation}

When MSA-DTI selects an alternative, the relative-margin rule gives
\begin{equation}
R_c(q_c^{*})
-
R_c(q^{\mathrm{ref}})
\leq
-\tau J_c(q^{\mathrm{ref}})
+
2\epsilon_c.
\label{eq:supp-expected-loss}
\end{equation}

Hence, the selected alternative has lower expected future loss than the
reference when
\begin{equation}
\tau J_c(q^{\mathrm{ref}})>2\epsilon_c.
\label{eq:supp-sufficient-condition}
\end{equation}
This is a sufficient condition and is not a guarantee for every target
center.

\subsection{Target-stage computational complexity}

Let $C_q^{\mathrm{grad}}$ be the cost of one target update for candidate
$q$, and let $C_q^{\mathrm{val}}$ be its validation cost. The target-stage
computational complexity is
\begin{equation}
O
\left(
|\mathcal{Q}|
+
\sum_{q\in\mathcal{Q}_{c}^{\mathrm{fea}}}
\left(
T C_q^{\mathrm{grad}}
+
C_q^{\mathrm{val}}
\right)
\right).
\label{eq:supp-target-complexity}
\end{equation}

The first term represents deployment filtering. The sum represents target
adaptation and validation of the feasible candidates. Offline candidate-bank
construction and relative-margin calibration are excluded.

\section{Additional Experimental and Diagnostic Results}
\label{supp:diagnostics}

\subsection{Target-side baseline settings}

The main manuscript keeps only the baseline information needed for the main
comparison. Table~\ref{tab:supp-baselines} gives the target-side settings
used in the main comparison.

\begin{suppTableBlock}{Target-side settings of MSA-DTI and the baselines in the main comparison.}{tab:supp-baselines}
\scriptsize
\renewcommand{\arraystretch}{1.08}
\setlength{\tabcolsep}{2.3pt}
\resizebox{\linewidth}{!}{%
\begin{tabular}{llllll}
\toprule
Method & Structure & Source initialization & Target update &
Adapted models & Deployment-limit check \\
\midrule
PT+FT & Fixed reference & Pooled-source model & SGD/MSE--50 & 1 &
Before output \\
MeDeT-based adaptation & Fixed reference & Meta-learned initialization &
SGD/MSE--50 & 1 & Before output \\
Random initialization & Fixed reference & Random & SGD/MSE--50 & 1 &
Before output \\
Few-shot NAS & Top-12 shortlist & Method-specific models & Adam/Huber--50 &
$\leq12$ & Before evaluation \\
Zero-shot NAS & Top-ranked model & Method-specific model & None & 0 &
Before evaluation \\
Zero-shot NAS + fine-tuning & Top-12 shortlist & Method-specific models &
Adam/Huber--50 & $\leq12$ & Before evaluation \\

H-Meta-NAS-based adaptation & Target architecture search &
Architecture-indexed meta-initializations & SGD/MSE--50 &
Multiple candidates & Before search evaluation \\

MSA-DTI & Candidate bank & Frozen source initializations & SGD/MSE--50 & 4--7 &
Before adaptation and output \\
\bottomrule
\end{tabular}}
\vspace{0.15em}

\footnotesize
PT+FT and MSA-DTI use the same source-trained reference and target
adaptation procedure. H-Meta-NAS-based adaptation uses the 50-update
budget-matched condition in the main comparison. It uses the same 66
architecture definitions, filters infeasible architectures before search,
uses source-only first-order meta-initializations, and selects the final
architecture using target validation data. Its settings are fixed without
held-out test data. All methods use the same target cases, deployment
limits, and architecture definitions.
\end{suppTableBlock}

\subsection{Optimizer-matched diagnostic comparison}

This comparison examines whether the performance ordering changes after the
target optimizer, task loss, and update budget are matched. All methods use
the same independent diagnostic target cases and the same deployment
limits.

\begin{suppTableBlock}{Full optimizer-matched comparison on the diagnostic target pool.}{tab:supp-matched}
\small
\renewcommand{\arraystretch}{1.08}
\begin{tabular}{lrrrr}
\toprule
Method & MSE$\downarrow$ & Worst-10\%$\downarrow$ &
CVaR90$\downarrow$ & Adapted models \\
\midrule
MSA-DTI & \textbf{0.003809} & \textbf{0.015094} &
\textbf{0.008696} & 6.10 \\
PT+FT & 0.004497 & 0.017144 & 0.011878 & 1.00 \\
Few-shot NAS & 0.006715 & 0.025246 & 0.019653 & 12.00 \\
Zero-shot NAS + fine-tuning & 0.007402 & 0.028099 & 0.020402 & 12.00 \\
Common shortlist, lowest validation loss &
0.006707 & 0.025240 & 0.019653 & 12.00 \\
Common shortlist, reference-based selection &
0.006819 & 0.025562 & 0.020547 & 12.00 \\
\bottomrule
\end{tabular}
\end{suppTableBlock}

MSA-DTI gives the lowest MSE, Worst-10\% error, and CVaR90 in this
comparison. The source initializations and selection rules still differ.
This comparison is therefore a robustness check rather than a complete
method equivalence test.

\subsection{Component variants}

The following diagnostic pool is separate from the source,
architecture-screening, margin-calibration, and held-out target pools.
It is used only for the component comparisons reported in this subsection.

The component variants are reported in Table~4 of the main text.

The historical source-training variant changes the complete procedure used
to construct the alternative source initializations and is therefore
interpreted as candidate-bank sensitivity rather than the isolated effect of
a single training choice.
The reference-only and two-reference-initialization variants show that
additional initializations of the reference architecture do not explain the
gain from cross-architecture selection.
Removing deployment filtering slightly lowers diagnostic MSE but reduces
limit satisfaction from 100\% to 96.25\%, confirming the role of filtering
in enforcing the hard limits.

\subsection{Selection-outcome details}

On the 80 independent diagnostic cases, the relative-margin rule selects
alternatives in 61 cases, including 49 beneficial and 12 harmful selections.
Direct lowest-validation-loss selection chooses alternatives in 69 cases,
including 55 beneficial and 14 harmful selections. The corresponding
all-case harmful-selection rates are 15.00\% and 17.50\%, respectively.


\paperfig{0.82\linewidth}{fig9.pdf}
{Selection outcomes on the 80-case diagnostic pool.}
{fig:supp-selection-diagnostic}

The relative-margin rule retains the reference more often. The confidence
intervals of the two harmful-selection rates overlap, so the observed
difference is not statistically confirmed.

\section{Timing, Hardware Profiling, and Deployment Analysis}
\label{supp:hardware}

\subsection{Target-side timing protocol}

All methods are implemented in Python and PyTorch and run on an NVIDIA RTX
3060 Laptop GPU. Target-side time includes model creation, initialization
loading, target adaptation, validation evaluation, and final selection. It
excludes offline preparation and test evaluation. Each reported time is the
mean over five runs, and the main manuscript reports the corresponding
standard deviation.

\subsection{Architecture-profiling protocol}

The five executable architectures are profiled with batch size one on an
Intel Core i7-12700H CPU and an NVIDIA RTX 3060 Laptop GPU for
$H\in\{1,4\}$. Each setting uses 100 warm-up inferences followed by five
repetitions of 1000 timed inferences. CUDA execution is synchronized before
and after each timed region. Median latency is the primary latency measure,
and the 95th percentile describes measurement variability. Serialized
parameter-file size is used as the measured model size. Peak allocated GPU
memory is also recorded for each architecture--horizon pair.

\begin{suppTableBlock}{Per-architecture batch-one hardware profile.}{tab:supp-hardware-profile}
\scriptsize
\renewcommand{\arraystretch}{1.05}
\setlength{\tabcolsep}{2.0pt}
\resizebox{\linewidth}{!}{%
\begin{tabular}{lrrrrrrr}
\toprule
Configuration & $H$ & Operations & Parameters & CPU (ms) & GPU (ms) &
State (KiB) & GPU peak (KiB) \\
\midrule
A1  & 1 & 155{,}712   & 77{,}921 & 0.097 & 0.148 & 306.5 & 1024 \\
A1  & 4 & 155{,}904   & 78{,}020 & 0.094 & 0.133 & 306.9 & 1024 \\
A6  & 1 & 157{,}760   & 78{,}977 & 0.093 & 0.168 & 311.0 & 1024 \\
A6  & 4 & 157{,}952   & 79{,}076 & 0.099 & 0.151 & 311.4 & 1024 \\
A13 & 1 & 159{,}808   & 80{,}033 & 0.128 & 0.237 & 315.5 & 1024 \\
A13 & 4 & 160{,}000   & 80{,}132 & 0.128 & 0.218 & 315.9 & 1024 \\
A55 & 1 & 377{,}888   & 2{,}081  & 1.793 & 3.999 & 10.3  & 54 \\
A55 & 4 & 377{,}984   & 2{,}132  & 1.793 & 4.004 & 10.4  & 1030 \\
A56/A57 & 1 & 1{,}050{,}688 & 5{,}697 & 1.980 & 4.019 & 24.4 & 60 \\
A56/A57 & 4 & 1{,}050{,}880 & 5{,}796 & 1.975 & 4.170 & 24.8 & 1037 \\
\bottomrule
\end{tabular}}
\vspace{0.15em}

\footnotesize
CPU and GPU values are median inference latency. A56 and A57 are the same
executable one-layer GRU-32 architecture. GPU peak is the maximum allocated
device memory recorded by the profiling protocol and is descriptive rather
than a configured memory limit.
\end{suppTableBlock}

\subsection{Profiling results}

Duplicate one-layer GRU-32 records are merged at the executable-architecture
level. Across the CPU/GPU measurements and $H\in\{1,4\}$, the Spearman
correlation between estimated operation count and median inference latency
ranges from 0.889 to 0.964. The correlation between parameter count and
serialized parameter-file size is 1.000 in every measured setting.
The MLP architectures use fewer estimated operations but more parameters
than the GRU architectures. The two deployment indicators therefore
describe different model properties, although this does not imply that both
limits are active under every tested deployment setting.
Peak GPU allocation is not monotonic in parameter count across the two
horizons, so the parameter-count limit is interpreted as a model-size
indicator rather than a direct device-memory limit. The measurements use
one CPU/GPU host and are not a cross-platform deployment benchmark.

\subsection{Single-limit feasible-space audit}

To determine which deployment limit is active in the frozen candidate bank,
we apply the operation-count and parameter-count limits separately to the
same seven candidates. This is a static complexity audit and does not repeat
source training, target adaptation, model selection, or test evaluation.

For both $H=1$ and $H=4$, all seven candidates satisfy the operation-count
limit under the tight, medium, and loose settings. Under the tight
parameter-count limit, the three MLP candidates are excluded and the
feasible bank is reduced from seven candidates to four, whereas the medium
and loose parameter-count limits retain all seven candidates. Consequently,
the joint feasible set contains four candidates under the tight setting and
seven under the medium and loose settings. The fixed reference candidate
satisfies both limits in all six horizon--tier combinations.

\clearpage
\null
\vspace{-1.5em}
\subsection{Expanded-bank dual-limit stress audit}

The preceding audit shows that the operation-count limit is inactive for
the frozen seven-candidate main bank under the main limit tiers.
To examine whether operation count and parameter count can independently
restrict candidate feasibility, we additionally construct an eight-candidate
stress bank by adding configuration 24 from the predefined architecture
space.

The stress thresholds are determined only from deterministic
model-complexity profiles, without using prediction-loss or target-performance
information. Across $H\in\{1,4\}$, the largest operation count in the
original seven-candidate bank is $1{,}050{,}880$, whereas the smallest
operation count of configuration 24 is $1{,}198{,}112$. Their midpoint gives
the stress operation-count limit $B^{F}=1{,}124{,}496$. Similarly, the
largest parameter count among the stress-bank candidates other than A13 is
$79{,}076$, whereas the smallest parameter count of A13 is $80{,}033$.
Their midpoint gives the stress parameter-count limit
$B^{P}=79{,}554.5$.

For both $H=1$ and $H=4$, configuration 24 violates only the
operation-count limit, A13 violates only the parameter-count limit, and the
fixed reference A57 satisfies both limits. Each single-limit feasible set
therefore retains seven of the eight candidates, whereas their joint
feasible set retains six. Thus, the joint feasible set is a strict subset
of each single-limit feasible set.

The corrected target-side stress evaluation uses the same 80 held-out cases,
and all outputs selected under the joint limits satisfy both limits.
This stress audit does not modify the frozen seven-candidate main bank,
the main deployment-limit tiers, or the headline experimental results.


\subsection{Selected-architecture analysis}

\paperfig{0.78\linewidth}{fig10.pdf}
{Selected-configuration complexity and performance. Marker area shows the
number of selected cases, and $h$ denotes harmful selections.}
{fig:supp-architecture-tradeoff}

The reference configuration is retained in 33 cases. Four alternative
configurations cover the remaining 47 cases, and their mean paired MSE
reductions range from 17.43\% to 30.50\%. All four alternatives have
positive mean reductions. The three harmful selections occur on two
alternative GRU configurations.

\section{Configuration Analysis}
\label{supp:configuration}

\paperfig{0.99\linewidth}{fig11.pdf}
{Configuration analysis of (a) the retained configuration count,
(b) the target update budget, and (c) the relative margin.}
{fig:supp-configuration}

Diagnostic MSE decreases from 0.004917 with one retained configuration to
0.004063 with four retained configurations and then remains unchanged. The
complete retained list is kept because it was frozen before this diagnostic
analysis. This result supports a compact candidate bank but does not
identify a globally optimal bank size.

Most of the adaptation gain appears within the first ten target updates.
The main experiment nevertheless uses 50 updates for every feasible
candidate to preserve a common fixed target budget. The relative-margin
analysis shows that 10\% is the smallest value satisfying all preset
calibration criteria.

\section{Source Sensitivity and External Transfer}
\label{supp:external}

\subsection{Source-center scale and source-training seeds}

\paperfig{0.98\linewidth}{fig12.pdf}
{Sensitivity to source-center scale, source-training seeds, and Alibaba
v2018 transfer.}
{fig:supp-source-sensitivity}

One Alibaba case below $-25\%$ is clipped in Fig.~\ref{fig:supp-source-sensitivity}
only for visualization; all original values are retained in the analysis.

The mean paired MSE reductions at source-center scales 10, 20, 30, 40, and
50 are 6.46\%, 4.92\%, 3.78\%, 3.41\%, and 4.14\%, respectively. The mean
gain remains positive at every tested source scale.

Across source-training seeds 2904, 2905, and 2906, the mean paired MSE
reductions are 10.52\%, 13.21\%, and 11.25\%. All three confidence
intervals remain above zero, although some individual cases have zero or
negative gains.
The screened architecture list is held fixed in this seed study. The result
therefore measures sensitivity to source-model training and initialization.

\subsection{Independent architecture-screening stability}

We separately repeat the complete 66-configuration diagnostic screening on
two new 20-center pools. Together with the original 20-center screening pool,
this gives three independently sampled screening sets. The A57 reference,
architecture space, initialization recipe, SGD/MSE--50 target recipe, and
retention rule are unchanged. A configuration is retained when it records at
least two check-oracle wins or at least two validation-selected positive
check wins. None of these pools overlaps the 20-center robustness pool used
below, and Test is never read.

\begin{suppTableBlock}{Independent candidate-screening stability and evaluation on centers 1160--1179.}{tab:supp-screening-stability}
\scriptsize
\renewcommand{\arraystretch}{1.08}
\setlength{\tabcolsep}{2.2pt}
\resizebox{\linewidth}{!}{%
\begin{tabular}{llllrrrr}
\toprule
Screen & Centers & Retained configurations & Size & Jaccard &
Gain (\%) & 95\% CI (\%) & Harmful (\%) \\
\midrule
Frozen S0 & 900--919 & A1, A6, A13, A55, A56, A57 & 6 & 1.000 & 0.295 & $[-0.646,1.564]$ & 2.50 \\
Independent S1 & 1120--1139 & reference fallback only & 0 & 0.000 & 0.000 & $[0.000,0.000]$ & 0.00 \\
Independent S2 & 1140--1159 & A55, A56, A57, A59 & 4 & 0.429 & 0.125 & $[-0.894,1.427]$ & 3.75 \\
\bottomrule
\end{tabular}}
\vspace{0.15em}

\footnotesize
Jaccard similarity is relative to the frozen S0 list. Size counts retained
configurations before adding the protected PT-A57 fallback. Gain is the mean
paired check-MSE reduction relative to PT-A57 on 80 robustness cases; intervals
use a 4000-repetition center-cluster bootstrap. This is a C1 screening
diagnostic and does not retrain or retune the frozen headline method.
\end{suppTableBlock}

Across S0--S2, retained-list sizes are 6, 0, and 4 (median 4, range 0--6).
A55, A56, and A57 are retained in two of three pools; A1, A6, A13, and A59
are each retained once. On the robustness pool, the S0 list selects PT-A57
in 72 cases, A55 in three, A56 in three, and the C1 A57 initialization in
two. S1 falls back to PT-A57 in all 80 cases. S2 selects PT-A57 in 71 cases,
A55 in three, A56 in three, C1-A57 in two, and A59 in one.

The rule therefore does not recover a stable or unique architecture set, and
neither nonempty list has a paired-gain interval excluding zero on the
independent pool. This negative result narrows the claim: the main results
evaluate a frozen, prespecified bank; they do not establish consistency of
the preceding candidate-discovery rule. It also supports keeping the
reference candidate outside the discovery decision and using it when no
alternative satisfies the retention rule.

\subsection{Alibaba external-domain transfer}

Alibaba Cluster Trace v2018 is split into 20 source machines, 20 calibration
machines, and 40 held-out machines.
The calibration and held-out pools contain 80 and 160 cases, respectively.
Models are retrained on the Alibaba source machines using their own
architectures and source-only normalization.
The architecture definitions, deployment filtering rule, target protocol,
and margin grid remain unchanged.

The workload observations are obtained from the production trace.
The tight, medium, and loose deployment-limit tiers remain controlled
model-complexity settings and are not direct measurements of machine-specific
latency, memory use, or energy use.
The Alibaba experiment therefore evaluates external workload-domain transfer
under the frozen deployment protocol.

No margin in the fixed grid satisfies all calibration criteria on Alibaba.
The previously frozen $\tau=0.10$ rule is therefore evaluated without
recalibration, and no Alibaba-specific calibrated margin is claimed.

On the 160 held-out cases, the adapted reference and MSA-DTI have mean MSE
values of 0.001337 and 0.001279, respectively, corresponding to a 4.29\%
aggregate reduction. The mean paired MSE reduction is 4.20\%, with a 95\%
confidence interval of $[2.03\%,6.70\%]$. MSA-DTI selects alternatives in
57 cases, including 47 beneficial and 10 harmful selections. The
harmful-selection rates are 6.25\% over all held-out cases and 17.54\%
conditional on alternative selection. The all-case harmful-selection rate
exceeds the original 5\% criterion, so the external evaluation shows a
positive average transfer gain but does not reproduce the original
reliability criterion.

\section{Synthetic Benchmark Generation and Center Heterogeneity}
\label{supp:synthetic-benchmark}

The controlled benchmark represents heterogeneous computing centers in a
common metric coordinate space.
Each center may observe only a subset of the metric positions.
The model input uses 12 value channels, 12 corresponding observation-mask
channels, and one time-gap channel, giving a fixed input dimension of 25.
Unavailable metric values are zero-filled and identified by the
corresponding mask channels.

Center types A--C are evaluation strata rather than model classes.
They differ in workload dynamics, sampling, and monitoring conditions.
Deployment-limit tiers are separate platform-side settings and are assigned
using the type-dependent probabilities reported in
Table~\ref{tab:supp-center-types}.
The table also gives the complete type-specific parameter ranges.
The default A:B:C ratio is 3:4:3, giving 6 Type-A, 8 Type-B, and 6 Type-C
centers in each 20-center source or held-out target pool.

\begin{suppTableBlock}
{Type-specific parameter ranges of the controlled synthetic benchmark.}
{tab:supp-center-types}
\scriptsize
\renewcommand{\arraystretch}{1.06}
\setlength{\tabcolsep}{3.2pt}
\begin{tabularx}{0.98\linewidth}{
    >{\raggedright\arraybackslash}p{0.28\linewidth}
    Z{0.21\linewidth}
    Z{0.21\linewidth}
    Z{0.21\linewidth}}
\toprule
Parameter & Type A & Type B & Type C \\
\midrule

Sampling interval $d_t$
& $\{1\}$
& $\{1,2\}$
& $\{2,3,4\}$ \\

Point-missing rate
& $[0.02,0.08]$
& $[0.10,0.25]$
& $[0.25,0.45]$ \\

Missing-block length
& $[5,15]$
& $[20,60]$
& $[60,120]$ \\

Observation-noise std.
& $[0.01,0.03]$
& $[0.03,0.06]$
& $[0.06,0.10]$ \\

Permanently unavailable metrics
& $0$--$2$
& $2$--$4$
& $4$--$6$ \\

Base workload logit
& $[-0.70,0.10]$
& $[-0.45,0.35]$
& $[-0.20,0.55]$ \\

Seasonal amplitude
& $[0.20,0.45]$
& $[0.30,0.60]$
& $[0.40,0.75]$ \\

AR coefficient
& $[0.55,0.80]$
& $[0.60,0.88]$
& $[0.65,0.92]$ \\

Process-noise std.
& $[0.04,0.08]$
& $[0.06,0.11]$
& $[0.08,0.14]$ \\

Burst rate per 1000 base steps
& $[1.0,2.5]$
& $[1.5,3.5]$
& $[2.0,4.5]$ \\

Burst amplitude
& $[0.25,0.55]$
& $[0.35,0.75]$
& $[0.45,0.95]$ \\

Drift probability
& $0.20$
& $0.40$
& $0.60$ \\

Drift strength
& $[0.10,0.25]$
& $[0.15,0.35]$
& $[0.20,0.45]$ \\

Limit-tier probabilities
(tight / medium / loose)
& $0.25/0.45/0.30$
& $0.35/0.45/0.20$
& $0.45/0.40/0.15$ \\

\bottomrule
\end{tabularx}
\end{suppTableBlock}

The three types share the same missingness mechanisms.
For every center, the generator may create 2--6 single-metric outage
blocks and 1--3 grouped outage blocks containing 2--4 metrics.
A short global monitoring outage occurs with probability 0.20.
The shared periodic base scales are 24, 96, and 672 base steps.

Together, these settings vary five aspects of the benchmark: workload
dynamics, sampling, monitored-variable availability, observation conditions,
and deployment-limit settings.

For each center, the generator first creates a latent workload process.
The workload contains periodic variation, autoregressive variation, and
random workload bursts.
When drift occurs, it is generated as either an abrupt change or a gradual
transition.

The latent workload is then mapped to multivariate monitoring signals.
Individual metrics may have different delays, gains, offsets, nonlinear
terms, and observation noise.
Center-specific sampling is applied before monitoring-schema and missingness
effects are introduced.

The benchmark generator and all data-role assignments are fixed before
architecture screening, margin calibration, diagnostics, and held-out
evaluation.

\end{document}
